\paragraph{消費者物価テイラールール（TR-CPI）}
\label{sec:chapter5_a_shock_policies_tr_cpi}
TR-CPIの式については
\ref{sec:chapter5_beta_shock_policies_tr_cpi}
を参照のこと。
\begin{enumerate}
    \item
    \label{itm:chapter5_a_shock_policies_tr_cpi_c_H_W}
        \(c_t^{H\to W}\)：
        総消費指数\(c_t^{H\to W}\)の図\ref{fig:chapter5_a_shock_graphs_c_H_W}を見ると、TR-CPIの軌道は消費者物価インフレ目標（IT-CPI）や消費者物価水準目標（CPLT）などと同様に、ショック直後の初期において大きなマイナス方向への大暴落を引き起こしていることがわかる。
        総消費水準目標（CLT）や名目総消費水準目標（NCLT）が落ち込みをほぼゼロに抑え込んでいるのとは対照的である。
        供給ショックに伴うインフレ圧力に対して、テイラールールに基づく機械的な金利の引き上げ（金融引き締め）で対応してしまうため、実体経済への強い下押し圧力が働き、消費の深刻な減少を招いている。
        消費の低下は厚生にマイナスに働くため、この前半における深刻な消費の落ち込みがTR-CPIの厚生を全体の中で下位クラスへと低迷させる大きな要因となっている。
        また深く落ち込んだ状態から長い時間をかけて定常状態へ回帰していく軌道を描いているため、消費増加による厚生の改善効果も得られていない。
    \item
    \label{itm:chapter5_a_shock_policies_tr_cpi_y_H}
        \(y_t^H\)：
        生産\(y_t^H\)の図\ref{fig:chapter5_a_shock_graphs_y_H}を見ると、TR-CPIの軌道は総消費指数\(c_t^{H\to W}\)と同様に、消費者物価インフレ目標（IT-CPI）や消費者物価水準目標（CPLT）と極めて類似した大暴落を見せていることがわかる。
        前述の通り、インフレ圧力に対する機械的な金融引き締めが実体経済をさらに圧迫するため、ショック直後の前半においてTR-CPIはマイナス方向へ極めて大きく落ち込んでいる。
        生産の低下は労働の減少を意味するため、この前半部分においては労働の非効用を減らして厚生にプラスに働く効果が他の政策よりも大きくなっている。
        しかしながら、前述した深刻な消費の大暴落によるマイナス効果や、後述する価格分散の増大によるマイナス効果を補うほどの改善には至っておらず、結果として全体的な厚生の評価を低迷させる結果となっている。
        また深く落ち込んだ状態から長い時間をかけて定常状態へ回帰する軌道を描いている。
    \item
    \label{itm:chapter5_a_shock_policies_tr_cpi_pi_H}
        \(\pi_t^H\)：
        グロス・インフレ率\(\pi_t^H\)の図\ref{fig:chapter5_a_shock_graphs_pi_H}を見ると、TR-CPIの軌道は消費者物価インフレ目標（IT-CPI）や消費者物価水準目標（CPLT）などと同様に、ショック直後の初期においてマイナス方向（デフレ）への下落を引き起こしていることがわかる。
        総消費水準目標（CLT）が初期に大きなプラス方向へのスパイクを記録しているのや、生産者物価インフレ目標（IT-PPI）が全期間を通じてゼロ近傍で極めて平坦な推移を維持しているのとは対照的である。
        テイラールールに従って消費者物価の変動に機械的な金融引き締めで対応してしまうため、実体経済を過度に冷やし込み、結果として国内生産財のインフレ率の落ち込みを招いている。この初期におけるインフレ率の変動が次項で述べる価格分散の蓄積をもたらし、結果として厚生を悪化させる一因となっている。
    \item
    \label{itm:chapter5_a_shock_policies_tr_cpi_Delta_t_minus_1_H}
        \(\Delta_{t-1}^H\)：
        価格分散\(\Delta_{t-1}^H\)の図\ref{fig:chapter5_a_shock_graphs_Delta_H}を見ると、TR-CPIの定常状態からの乖離は、生産者物価水準目標（PPLT）や生産者物価インフレ目標（IT-PPI）が全期間を通じてほぼゼロの水準を維持していることと比較して、明確な増大を見せていることがわかる。
        これは前述のインフレ率の初期における下落が価格分散として蓄積した結果であり、この生産効率の悪化がTR-CPIの厚生を押し下げる一因となっている。
        しかしながら、供給ショック（\(a\)ショック）においては価格分散の増大によるマイナス効果以上に、実体経済（消費や生産）の大暴落によるマイナス効果が極めて絶大である。そのため、TR-CPIの厚生が総消費水準目標（CLT）などに比べて低迷している主たる要因は、価格分散の蓄積よりもむしろ初期の深刻な不況を防げなかった点にあると言える。
\end{enumerate}